\documentclass[11pt]{article}

\usepackage[margin=1in]{geometry}
\usepackage[T1]{fontenc}
\usepackage[utf8]{inputenc}
\usepackage{lmodern}
\usepackage{microtype}
\usepackage{amsmath,amssymb,amsthm,mathtools}
\usepackage{booktabs}
\usepackage[numbers,sort&compress]{natbib}
\usepackage{hyperref}

\hypersetup{
  colorlinks=true,
  citecolor=blue,
  linkcolor=blue,
  urlcolor=blue
}

\newcommand{\code}[1]{\texttt{#1}}
\newcommand{\Z}{\mathbb{Z}}
\newcommand{\bits}[1]{\operatorname{bits}(#1)}
\newcommand{\cost}[1]{\operatorname{cost}(#1)}
\setlength{\emergencystretch}{3em}

\title{Certified Polynomial-Complexity Hermite and Smith Paths in Lean 4:\\
Explicit Inverses and a Kannan--Bachem Schedule}
\author{Junye Ji\\University of Washington}
\date{Research version 0.1.0, July 22, 2026}

\begin{document}

\maketitle

\begin{abstract}
The classical Kannan--Bachem result separates the mathematical correctness of
integer Hermite and Smith normal forms from the polynomial control of their
intermediate values.  We present an executable Lean 4 development of this
separation for nonsingular square integer matrices.  A principal-minor HNF
kernel is preceded by a deterministic certified row permutation.  Its minor
selection uses a division-free Bird determinant whose executable value,
intermediate widths, and schoolbook arithmetic cost share one verified
implementation.  A total SNF path then alternates one-column left HNF and
transposed prepared HNF, clears divisible columns, and performs the
Kannan--Bachem Step 7 injection when the pivot fails to divide the lower
block.  Every continuation strictly decreases the binary size of the pivot.
Both algorithms return the project's existing strong result types, including
forward transforms, explicit inverses, both inverse identities, and canonical
normal-form proofs.  Kernel theorems separately establish semantic
correctness, ring-operation bounds, coefficient-width bounds for every
transform and inverse, and a composed binary bit-operation bound.  A pinned
artifact audits 263 theorem roots, freezes 572 public declarations, and
records three deterministic native cases.  Native timing remains regression
evidence and is not used to justify the formal complexity claim.
\end{abstract}

\section{Introduction}

An integer normal-form routine can be correct and still be a poor algorithm.
Elementary row and column operations preserve the relevant equivalence, but
uncontrolled choices can produce enormous intermediate coefficients.  The
Kannan--Bachem analysis was the first to bound both the number of algebraic
operations and the binary length of intermediates by polynomials in the input
size~\citep[pp. 499--505]{kannanBachem1979}.  A mechanized treatment must expose more
than the final normal-form predicate: it must connect the executed branches,
transformation matrices, termination measure, coefficient envelopes, and a
concrete integer cost model.

This work supplies that connection through a nonsingular-square research
seam in the project.  It is available through the additive facade
\begin{center}
\code{NormalForms.Research.KannanBachem}.
\end{center}
The stable core facade does not re-export it.  The implementation uses Lean 4
and mathlib~\citep{deMouraUllrich2021Lean4,mathlib2020} and returns the same strong HNF and SNF
types as the core algorithms.  Consequently the research path can be compared
with the canonical core result without inventing a second mathematical API.

The main contributions are:
\begin{itemize}
  \item a primitive-traced principal-minor HNF schedule with directly
        accumulated forward and inverse row transformations;
  \item deterministic row preparation for every nonsingular square input,
        driven by a costed, division-free determinant evaluator;
  \item a row-oriented adaptation of Smith Steps 4--7, including
        the nondivisible lower-block injection and a strict pivot-width
        termination proof;
  \item explicit polynomial bounds for algebraic operations and all
        intermediate coefficients in the result and four transformations;
  \item a binary schoolbook cost recurrence whose value projection follows
        the executed arithmetic and whose closed endpoint depends only on
        dimension and original input width;
  \item a reproducible, independently versioned paper, artifact, axiom audit,
        API freeze, native corpus, and raw benchmark profile
        \citep{ji2026normalforms}.
\end{itemize}

The scope is exact.  Research version 0.1.0 proves these results for
nonsingular matrices in \(\Z^{n\times n}\).  It does not claim that this
particular schedule handles rectangular or rank-deficient inputs.  The stable
core already handles those inputs; a future Kannan--Bachem wrapper would need
its own rank decomposition and complexity proof.

\section{Interfaces and certificate discipline}

For an input matrix \(A\), the HNF result stores \(U,H\), an explicit matrix
\(U^{-1}\), and proofs
\[
  UA=H,\qquad U^{-1}U=I,\qquad UU^{-1}=I.
\]
The SNF result analogously stores \(U,S,V,U^{-1},V^{-1}\) and proves
\[
  UAV=S,qquad U^{-1}U=UU^{-1}=I,
  \qquad V^{-1}V=VV^{-1}=I.
\]
The inverse matrices are data, not objects recovered afterward through a
determinant theorem, a matrix inverse operator, or choice.  Every primitive
row or column step carries its inverse, and composition multiplies forward
and backward transforms in the correct opposite orders.

This discipline matters for complexity.  A theorem about \(S\) alone says
nothing about the size or cost of the equivalence witness.  The profile
therefore bounds five observable matrices for SNF:
\[
  S,\quad U,\quad U^{-1},\quad V,\quad V^{-1}.
\]
The Lean structure \code{SmithCoefficientProfile} records their actual bit
widths.  The endpoint theorem bounds each field using the original dimension
and input coefficient width.

Operation logs and final certificates have different roles.  A
\code{RowTrace} is an executable list of swaps, additions, unit scalings, and
two-by-two B\'ezout pair steps.  Replaying the trace explains the algorithm's
primitive schedule.  The accumulated certificate is the stable public
result.  Theorems connect replay to the result matrix and connect both
accumulators to the explicit inverse identities.

\section{Prepared principal-minor HNF}

Kannan--Bachem state their HNF schedule in a column convention.  This section
uses its transposed row form: left row transformations replace the original
right column transformations, while the principal-minor and coefficient
arguments are transported unchanged \citep[Algorithm HNF and Theorem
2]{kannanBachem1979}.

\subsection{The ready kernel}

The square HNF kernel processes leading principal blocks.  Its explicit
precondition \code{PrincipalReady A} states that every canonical leading
principal submatrix has nonzero determinant.  Under this invariant, each
stage performs a bounded extended-gcd row combination, normalizes the pivot,
and reduces entries above it.  Arithmetic events align one-for-one with the
primitive trace, so the operation count is derived from the executed path
rather than reconstructed from the final matrix.

For dimension \(n\), the trace length is at most \(n^3\), a concrete refinement
of the polynomial operation bound in the cited HNF analysis.  The corresponding
\code{RingOperationCount} separates additions, multiplications, extended-gcd
calls, and normalizations.  This separation avoids assigning unit cost to an
integer gcd when the later theorem needs its binary cost.

\subsection{Removing the readiness assumption}

Nonsingularity does not imply that the input's current leading principal
minors are all nonzero.  The function \code{Principal.prepare} constructs a
row permutation recursively.  At each level it inspects complementary minors
formed by deleting a candidate row and the last column, chooses the first
nonzero candidate, places that row last, and recurses on the retained block.
Nonsingularity proves that a candidate exists.

The preparation object stores the permutation, permuted matrix, permutation
matrix, its explicit inverse certificate, the transformation equation, and a
proof that the output is ready.  The entry
\code{preparedPrincipalHermiteNormalForm} composes preparation with the ready
kernel and proves equality with the core canonical HNF.  Thus callers provide
only \(\det A\ne0\); the leading-minor obligation does not leak through the
public research entry point.

\subsection{A charged determinant selector}

Minor selection must not hide determinant work.  Preparation therefore uses
a cached dense-vector implementation of Bird's division-free determinant
recurrence~\citep[p. 1072]{bird2011determinants}.  Every scalar addition and
multiplication is performed by the verified binary sign-magnitude reference
model.  The implementation returns both a value and an exact modeled cost.

The proof stack establishes:
\begin{enumerate}
  \item decoding a dense-vector stage equals the algebraic matrix recurrence;
  \item the recurrence has closed height and bit-length envelopes;
  \item the costed scalar and matrix iterations satisfy a closed schoolbook
        budget;
  \item the recurrence equals the determinant for arbitrary dimension;
  \item the first-nonzero scan returns exactly the first minor whose
        determinant is nonzero and is bounded across all recursive levels.
\end{enumerate}
The theorem \code{evaluateWithCost\_value\_eq\_det} is the semantic bridge.
The end-to-end prepared-HNF cost adds the actual recursive determinant scans
to the ready-kernel arithmetic.

\section{Total Smith Steps 4--7}

The presentation below is the row-HNF transpose of Kannan--Bachem's
column-oriented Smith schedule \citep[Algorithm SNF, Steps 4--7,
p.~505]{kannanBachem1979}.

\subsection{One pass}

The active block is nonsingular and nonempty.  Step 4 applies a dedicated
one-column left-HNF reducer to the entire block.  It uses exactly one bounded
extended-gcd pair step per lower row and one final normalization.  Step 5
transposes the result, runs prepared principal HNF, and transports the row
multiplier back as a right transformation.  The resulting
\code{TwoSidedCertificate} already contains both explicit inverses.

After a pass, the pivot is nonzero and normalized and the first row is
cleared.  Entries below the pivot fall into two cases.  If one is not
divisible by the pivot, the next pass replaces the pivot by a proper common
divisor.  If all are divisible, one certified shear clears the complete first
column at once.  These are the repeated Steps 4--6 of the cited schedule.

\subsection{Step 7 and termination}

Once the first row and column are clear, the pivot may still fail to divide an
entry in the lower-right block.  The search function returns the first such
position.  Step 7 adds its column into the pivot column, thereby placing the
offending entry below the pivot without changing the pivot itself.  The next
pass again produces a proper common divisor, matching the cited Step 7 after
the row/column transpose.

The total function \code{stabilizeFrom} follows exactly these branches.  Its
well-founded measure is
\[
  \bits{A_{00}} = \operatorname{Nat.size}(|A_{00}|).
\]
Each recursive branch proves strict decrease of this measure.  No fuel,
factorization, or fallback to the core Smith algorithm appears in the primary
entry.  A separate fuelled function remains useful for diagnostics.

At termination the pivot is normalized and nonzero, both borders are clear,
and the pivot divides the full lower-right block.  Structural recursion then
runs the same algorithm on that lower-right block.  Forward and inverse lower
certificates are lifted through a leading identity block and composed with
the stabilized certificate.  The resulting \code{smith A hdet} directly has
type \code{SNFResultFin A}.

The semantic endpoints prove the transformation equation, inverse recovery,
the Smith predicate, and equality with the stable core canonical Smith
matrix.  Equality is asserted only for the canonical result; nonunique
transformation matrices are intentionally not compared.

\section{Four levels of complexity evidence}

The project uses four acceptance levels that must not be conflated.

\paragraph{Semantic correctness.}
The algorithms return the required normal form and strong transformation
certificate.  Termination and canonical equality are kernel checked.

\paragraph{Ring-operation count.}
Exact counters mirror the recursion and separately record additions,
multiplications, extended-gcd calls, normalizations, and exact divisions.  A
dimension- and input-width polynomial bounds the total.  Searches,
comparisons, allocation, and data movement are not mislabeled as ring
operations.  This discharges the operation-count obligation corresponding to
Kannan--Bachem's HNF Theorem 2 and SNF Theorem 4
\citep[pp. 503--505]{kannanBachem1979}.

\paragraph{Coefficient bit length.}
Bounds cover intermediate active matrices, arithmetic operands, forward
prefixes, inverse prefixes, clear and injection transformations, and every
certificate composition.  Determinant invariance keeps lower-right recursive
levels within a uniform original-input envelope.  This is the formal
coefficient-control obligation represented by their HNF Lemmas 1--3 and SNF
Theorem 5 \citep[pp. 502--505]{kannanBachem1979}.

\paragraph{Binary bit complexity.}
The scalar model uses canonical sign-magnitude words, carry/borrow addition,
schoolbook multiplication, shared Euclidean long division, and costed
extended gcd; see standard distinctions between algebraic and bit
cost~\citep[Chapter 1]{brentZimmermann2010}.  The matrix layer charges the actual scalar
reference operations rather than assuming native \code{Int} takes unit time.

For dense certificate multiplication, \code{matrixEntryWithCost} evaluates a
dot product through costed sign-magnitude multiplication and addition.
\code{matrixProductWithCost\_value} proves each decoded entry equals ordinary
\code{Matrix.mul}.  Composing two strong certificates charges four such
products in their true orders:
\[
  U_2U_1,\quad U_1^{-1}U_2^{-1},\quad
  V_1V_2,\quad V_2^{-1}V_1^{-1}.
\]

The exact recursive function \code{smithBitOperationCost} follows all
stabilization branches and every lower-right call.  Its closed endpoint is
\[
\begin{split}
 &\code{smithBitOperationCost A hdet}\\
 &\qquad\le
 \code{smithPolynomialBitOperationBound n (matrixBitLength A)}.
\end{split}
\]
The right-hand side depends only on the original dimension and maximum input
coefficient width.  It is intentionally conservative: the theorem's purpose
is a complete compositional polynomial bound, not a tight performance model.

\section{Theorem and artifact map}

Table~\ref{tab:endpoints} names the principal public endpoints.  The facade
freeze checks these and their supporting interfaces at compile time.

\begin{table}[ht]
\centering
\small
\begin{tabular}{p{0.25\linewidth}p{0.66\linewidth}}
\toprule
Layer & Public endpoint or evidence \\
\midrule
Prepared HNF semantics &
\code{preparedPrincipalHermiteNormalForm\_eq\_reference} \\
HNF operation count &
\code{principalRingOperations\_total\_le} \\
HNF coefficient width &
\code{principalIntermediateMatrixBitLength\_le\_polynomial\_of\_ready} \\
Prepared HNF bit cost &
\code{preparedPrincipalHNFBitOperationCost\_le} \\
SNF semantics &
\code{smith\_sound} and \code{smith\_eq\_reference} \\
SNF operation count &
\code{smithOperationCount\_total\_le\_polynomial} \\
SNF coefficient width &
\code{smith\_coefficientProfile\_le\_polynomial} \\
SNF bit cost &
\code{smithBitOperationCost\_le\_polynomial} \\
Trust audit & 263 roots; exact allowlist
\code{propext}, \code{Quot.sound}, \code{Classical.choice} \\
Surface freeze & 572 compile-time \code{\#check} declarations \\
\bottomrule
\end{tabular}
\caption{Research version 0.1.0 endpoints.}
\label{tab:endpoints}
\end{table}

The allowlisted logical axioms arise in the mathlib proof layer.  The audit
rejects project-defined axioms, \code{sorryAx}, and native/compiler trust.
Native execution is never used to manufacture a theorem.

\section{Executable evaluation}

The fixed corpus contains three branch-sensitive cases.  Table
\ref{tab:cases} records deterministic values; these are checked on every
artifact run before timing is considered.

\begin{table}[ht]
\centering
\small
\begin{tabular}{lrrrrrr}
\toprule
Case & \(n\) & input bits & ring ops & XGCD & passes/inj. & modeled cost \\
\midrule
prepared HNF & 3 & 5 & 93 & 3 & 0/0 & 242765386189266 \\
Smith repeat & 2 & 2 & 48 & 2 & 1/0 & 1130326542472 \\
Smith injection & 2 & 2 & 96 & 4 & 2/1 & 3012432681653 \\
\bottomrule
\end{tabular}
\caption{Deterministic executable telemetry.  ``inj.'' means Step 7
injections.  Modeled costs are not nanoseconds.}
\label{tab:cases}
\end{table}

The HNF input is
\[
\begin{pmatrix}2&3&5\\4&7&11\\6&13&17\end{pmatrix},
\qquad
H=\begin{pmatrix}2&0&0\\0&1&1\\0&0&2\end{pmatrix}.
\]
The two Smith inputs are
\[
\begin{pmatrix}2&1\\0&2\end{pmatrix}
  \longmapsto \operatorname{diag}(1,4),
\qquad
\begin{pmatrix}2&0\\0&3\end{pmatrix}
  \longmapsto \operatorname{diag}(1,6).
\]
The second Smith case forces Step 7 because the first pivot does not divide
the lower-right entry.

The raw benchmark performs one warmup and seven measurements for prepared
HNF, repeated-pass SNF, injection SNF, and the complete three-case corpus.  It
stores every observation rather than only an aggregate, computes median and
inclusive-quartile IQR, and records wall time, user/system time, maximum RSS,
source revision and patch digest, binary SHA256, Lean version, CPU, memory,
kernel, libc, and Python version.  Absolute wall time is compared only when
the hardware fingerprint matches.  The committed JSON and CSV are the source
of truth; a plotted or prose summary cannot override them.

On the captured host, median wall times are \(4.20\) seconds for prepared HNF,
\(0.06\) seconds for repeated-pass SNF, \(2.25\) seconds for injection SNF,
and \(6.38\) seconds for the complete corpus.  Median maximum resident sets
are respectively 71,980, 72,760, 73,104, and 72,348 KiB.  These numbers
describe this reproducibility run only; the JSON retains the raw observations
and IQR needed to interpret them.

The runtime checks recompute transformation equations, inverse recovery,
expected canonical diagonals, and actual-cost inequalities.  These checks are
useful regression witnesses but are weaker than the imported proofs.  In
particular, an executable returning success does not add a theorem to the
kernel.

\section{Reproducibility and trust boundary}

The host gate builds the research facade, its public API and execution tests,
the core facade isolation test, and the native benchmark.  It invokes Lean
directly on the execution regression file so that every \code{\#guard} is
re-elaborated even when Lake has cached an object.  It then runs the complete
test suite, all native cases, the exact axiom audit, a fresh benchmark, the
paper build, metadata validation, and \code{git diff --check}.

The pinned container fixes Debian 13.1, the Lean 4.32.0 binary archive and its
SHA256, and the committed Lake manifest.  The profile has no FLINT dependency.
FLINT is an optional untrusted certificate generator elsewhere in the
project, not a premise of any theorem in this development.

The trust base is Lean's kernel, the Lean implementation and parser used to
check source, and the ordinary assumptions represented by the exact axiom
report.  The strict research claim does not rely on \code{native\_decide},
\code{Lean.ofReduceBool}, \code{Lean.ofReduceNat}, external C arithmetic, the
benchmark harness, or the operating system's timer.

\section{Limitations and future work}

The present result should not be generalized beyond its statement.
\begin{itemize}
  \item Inputs are nonsingular square integer matrices.  A rectangular and
        rank-deficient wrapper needs a verified fraction-free rank
        decomposition and a proof that its transformations and coefficient
        bounds compose with this kernel.
  \item The modeled cost uses transparent schoolbook arithmetic.  It does not
        claim to model Lean's optimized native arbitrary-precision integers.
  \item Searches and allocation are excluded from the bit-operation count.
        They remain visible in native wall time and RSS.
  \item Bounds are explicit and polynomial but deliberately loose.  Tighter
        exponents and constants are separate optimization work.
  \item This profile does not make the Kannan--Bachem path the core default.
        Changing the default backend would require independent compatibility
        and performance evidence.
\end{itemize}

Related formalizations establish Smith normal-form correctness in other proof
assistants~\citep{divasonThiemann2022snf}.  The distinguishing goal here is the
end-to-end connection between a row-oriented executable adaptation, explicit
inverse transformations, intermediate coefficient control, and a concrete
binary arithmetic cost.

\section{Conclusion}

Research version 0.1.0 closes four distinct proof obligations for
nonsingular-square integer HNF and SNF: semantics, algebraic operation count,
coefficient growth, and bit complexity.  The key engineering decision is to
make every source of hidden work explicit.  Row preparation charges its
determinant scans; extended gcd shares the verified primitive model;
termination exposes a decreasing pivot width; and certificate composition
charges the forward and inverse matrix products.  The result is an
independently reproducible research line whose claims are precise enough to
audit and whose limitations do not weaken the stable general-purpose core.

\bibliographystyle{plainnat}
\bibliography{../references}

\end{document}
